\section{Background on Query Discovery}
\label{sec:problem-setting}
%\todo[inline]{R2: The explanation of the query model is good. However, it
%could be supplemented by different pattern-matching semantics
%(contiguous, skip-till-next, skip-till-any) and an explanation of how the
%window parameter can be learned.
%
%A: We will add a statement on how to incorporate further query semantics and
%how to determine the window size.}
Below, we introduce preliminary notions to study
the problem of event query discovery.

%\subsection{Event Streams and Queries}
%\label{sec:streams}

%We adopt some standard notation for sequences, as follows. For a set
%${Y}$, the set of all sequences over ${Y}$ is ${Y}^*$.
%For a finite sequence $z=\langle y_1,\ldots, y_n \rangle \in {Y}^*$, we
%write $z[i] = y_i$ for the $i$-th element.

\sstitle{Event streams}
We based our work on the notion of a multivariate event
stream. Such a stream has an \emph{event schema}, captured
by a tuple of attributes $\mathcal{A}= (A_1, \ldots,
A_n)$, where $\dom(A_i)$ is the finite domain of attribute
$A_i$, $1\leq i\leq n$. All events have the same schema
and we assume that attribute domains are disjoint,
$\cap_{i=1}^n \dom(A_i) =
\emptyset$. As a short-hand, we write $\Gamma =
\bigcup_{i=1}^n{\dom(A_i)}$ for the alphabet defined by
all possible values of attributes. Moreover, an
\emph{event} is an instantiation of the event schema,
captured as a tuple $e=(a_1, \ldots, a_n)$ with $a_i\in
\dom(A_i)$ for $1\leq
i\leq n$. By $e.A_i$, we denote the value of attribute
$A_i$ of event $e$.

An \emph{event stream} is a finite
sequence of events $S=\langle e_1, \ldots, e_l
\rangle$, i.e., a finite subsequence of a
potentially unbounded sequence of events. We write $|S|$
for the length of the stream, and $S[i]$ for its $i$-th
event. A \emph{stream database} is a finite set of
streams $D=\{S_1,\ldots, S_d\}$.

\sstitle{Event queries}
We adopt a query model that captures sequences of events,
which are correlated by predicates over their attribute
values~\cite{DBLP:conf/sigmod/ZhangDI14}. We neglect the
definition of an explicit time window, as it may be
derived from the maximal time difference between events in
one of the streams. That is, taking the time between the start and end of
each stream, the maximal time span over all streams is an upper bound for the
time interval, in which a situation of interest is expected to occur.
Moreover, we
follow~\cite{DBLP:conf/btw/Kleest-Meissner23} and use a
linearized representation of event queries.
Let $\mathcal{A}= (A_1, \ldots,
A_n)$ be an event schema, $\mathcal{X}$ be a finite set of
variables
such that $\mathcal{X}\cap
\Gamma=\emptyset$, and $\_ \notin \mathcal{X}\cup
\Gamma$ be a placeholder symbol. Then, we define a
\emph{query
term} $t=(u_{1},\ldots,u_{n})$ as an $n$-tuple comprising
attribute
values, variables, and placeholders:
$u_{j}\in \dom(A_j) \cup \mathcal{X} \cup \{ \_ \}$,
$1\leq j\leq n$, while we prohibit terms containing only
placeholders, i.e., $u_{j}\in
\dom(A_j) \cup \mathcal{X}$ for at least one $1\leq j\leq n$.
By $\varepsilon$, we denote the empty query term that
contains only
placeholders, i.e., %$\varepsilon=(\underbrace{\_,\dots,\_}_{n})$
$\varepsilon=(\_,\dots,\_)$.
Based thereon, an \emph{event
query} is a finite sequence of query terms, $q=\langle
t_1,\ldots,t_k \rangle$,
with $q[i]$ denoting the $i$-th term. In a query, variables need to occur in at
least two query terms for the same attribute, i.e., for any term
$q[i]=(u_{1},\ldots,u_{n})$ with $u_j\in \mathcal{X}$,
$1\leq i\leq k$ and $1\leq j\leq n$, there exists another query term
$q[p]=(u'_{1},\ldots,u'_{n})$ with $u_j = u'_j$,
$1\leq p\leq k$ and $p\neq i$. The empty query is
$\langle\varepsilon\rangle$, i.e., the query containing
the empty query term. By $\mathcal{Q}$, we denote the
universe of all queries.

From an intuitive point of view, a query term describes a
single event that shall be matched by defining a specific
value, variable, or placeholder per attribute. While a
placeholder denotes the absence of any constraint, the use
of variables captures the correlation of events as
matching events need to have equivalent values for the
respective attributes.


%\item $\tau\in \mathbb{N}$ is a window with $\tau \leq k$.
%\end{itemize}
%Given an event
%schema $\mathcal{A} =
%(A_1, \ldots, A_n)$ and a finite set of variables $\mathcal{X}$,
%an event query
%%is a tuple $q=(T,\tau)$, where
%%\begin{itemize}[nosep]
%%\item
%is a finite sequence of query terms $q=\langle t_1,\ldots,t_k \rangle$,
%with $q[i]$ denoting the $i$-th term. Each term
%$t_i=(u_{1},\ldots,u_{n})$, $1\leq i\leq k$, is an $n$-tuple built of
%attribute
%values, variables, and the dedicated
%symbol $\varepsilon \notin \Gamma \cup \mathcal{X}$ to capture the absence of
%a
%constraint,
%i.e.,
%$u_{j}\in \dom(A_j) \cup
%\mathcal{X} \cup \{\varepsilon\}$, $1\leq j\leq n$.
%%\item $\tau\in \mathbb{N}$ is a window with $\tau \leq k$.
%%\end{itemize}

%Note that variables that occur only once in a query are
%semantically
%equivalent
%to placeholders and, as such, will be largely neglected in the remainder.
%Also,
%only the elements of $\mathcal{X}$ that occur at least twice in a query are
%referred to as variables.
%Adopting this terminology,

Moreover, we distinguish three classes of queries: A
\emph{type query} contains only query terms defining
attribute values (and
placeholders). A \emph{pattern query} contains only query
terms defining variables (and
placeholders). All other queries are referred to as
\emph{mixed queries}.

% only
%variables, or a combination of both.
%These are denoted as type queries, pattern queries, and mixed queries,
%respectively.
% To improve readability, the non specified attributes are left empty in the query
%terms.

\begin{table}[t]
	\caption{Overview of notations.}
	\label{tab:notations}
	\vspace{-.5em}
	\footnotesize
	\begin{tabularx} {\textwidth}{l
	 @{\hspace{.5em}} p{10cm}}
	\toprule
	Notation & Explanation \\
	\midrule
	$\mathcal{A}=(A_1, \ldots, A_n )$ & An event
	schema, a
	tuple of attributes\\
	$\Gamma  = \bigcup_{i=1}^n{\dom(A_i)}$ & The stream alphabet\\
	$e= (a_1, \ldots, a_n)$ & An event, a tuple
	of attribute values\\
	$e.A_i$ & The value $a_i$ of attribute $A_i$ of
	event $e$\\
	%$e.t$ & The occurrence time of event $E$\\
	$S=\langle e_1, \ldots, e_l \rangle$ & An event
	stream, a finite sequence of events\\
	$S[i]$ & The $i$-th event of the event stream $S$\\
	$D=\{S_1,\ldots, S_d \}$ & A stream database, a
	set of event streams\\
%	$\Gamma_D$ & The supported stream alphabet of stream database $D$\\
	\midrule
	$\mathcal{X}$ & A finite set of query variables, $\mathcal{X} \cap
	\Gamma=\emptyset$\\
	$\_$ & A placeholder symbol, $\_ \notin \mathcal{X}\cup
	\Gamma$\\
	$t=(u_{1},\ldots,u_{n})$& A query term, an
	$n$-tuple of attribute values, variables, and placeholders\\
%	with $u_{j}\in
%	\dom(A_j) \cup \mathcal{X} \cup \{\_ \}$, $1\leq j\leq n$, such that not
%	all its components are placeholders\\
	$\varepsilon=(\_,\ldots,\_)$& The empty query term, an
$n$-tuple of placeholders\\
	$q=\langle
	t_1, \ldots, t_k \rangle$ & An event query, a finite sequence of
	query terms \\
	$q[i]$ & The $i$-th term of the query $q$\\
	$\langle \varepsilon\rangle$ & The empty query containing only the empty
	query term \\
	\midrule
	$S \models q$ & The stream $S$ supports the
	query $q$, i.e., there exists a match\\
	$M_+(q,D)$ & The set of streams in $D$ that match the query $q$\\
%	$M_-(q,D)$ & The set of streams in $D$ that do not match the query $q$\\
	supp$(q,D)$ & The support of query $q$ in stream database $D$\\
	$\theta$ & The support threshold, a real number in the interval $(0,1]$\\
	$s(|D|,\theta)$ & The stream threshold, the min number of streams
	to match to fulfill the support threshold\\
	$\Gamma(D,\theta)$ & The frequent stream alphabet\\
	$Q_{D}^{\theta}$ & The set of all matching queries\\
	%$q$ in stream $S$\\
	\bottomrule
	\end{tabularx}
	\vspace{-1em}
\end{table}

We define the semantics of a query $q=\langle
t_1,\ldots,t_k \rangle$ over a stream $S=\langle e_1,
\ldots , e_l\rangle$ as follows: A \emph{match} of $q$ in
$S$ is an injective mapping
$m: \{1,\ldots,k\}\rightarrow \{1, \ldots, l\}$,
such that:
\begin{itemize}[nosep,left=.2em .. 1.5em]
	\item for each query term $q[i]=(u_{1},\ldots,u_{n})$, $1\leq i\leq k$, the
	mapped event $S[m(i)]=(a_1,\ldots, a_n)$ has the
	required attribute
	values, $u_j = a_j$ or $u_j\in \mathcal{X}\cup \{\_ \}$
	for $1\leq 	j\leq n$;
	\item variables are bound to the same attribute values, i.e., for query
	terms $q[i]=(u_1,\ldots, u_n)$ and $q[p]=(u'_1,\ldots, u'_n)$, $1\leq
	i,p\leq k$, it holds that $u_j\in \mathcal{X}$, $1\leq j\leq n$, with
	$u_j=u'_j$ implies that $S[m(i)].A_j = S[m(p)].A_j$; and
 	\item the order of events in the stream is preserved,
 	i.e., $m(i)<m(p)$ for $1\leq
	i<p\leq k$.
\end{itemize}
Note that the set of matches as induced by the above definition corresponds
to a greedy event selection
policy (also known as unconstrained or skip-till-any-match). It does not
limit how often an event can participate in matches. However, in the
remainder, we do not explicitly refer to the set of matches, but consider
only the existence of a match per stream. Specifically, we say that
$q$ matches $S$ and write $S \models q$, if there exists a
match for query $q$ in stream $S$. \autoref{tab:notations}
provides an overview of the introduced notations.

\sstitle{Support} To reason about the relation of a query
$q$ and a stream database $D$, we
define $M_+(q,D)$ as the
set of matching streams, i.e., $M_+(q,D) = \{S \in D \mid
S \models q\}$.
The \emph{support} supp$(q,D)$ of a query is then given as
the fraction of streams in $D$ that match the query,
supp$(q,D)= \frac{|M_+(q,D)|}{|D|}$.
A \emph{support threshold} $\theta$ is a real number in
the interval $(0,1]$.
A query $q$ is \emph{frequent} in $D$, if supp$(q,D) \geq
\theta$.
From the aforementioned notions, we derive the
\emph{stream threshold} $s(|D|,\theta)$ as the minimum
number of streams
that must match to a query in order for the query to be
considered frequent, i.e., $s(|D|,\theta) = \lceil |D| \cdot
\theta \rceil$.
Based thereon, we also derive the frequent stream alphabet, defined by all
attribute values that appear in at least the number of streams as
determined by the stream threshold, i.e., $\Gamma(D,\theta) = \{a \in
\Gamma \mid |\{S \in D \mid \exists \ j\in\{1,\dots,|S|\},
i\in\{1,\dots,n\}: S[j].A_i = a\}| \geq s(|D|,\theta)\}$.
Furthermore, as a short-hand,
we define the set of all matching queries as $Q_D^{\theta}
= \{q \mid \text{supp}(q,D) \geq \theta\}$.

For a given stream database, it is clear that support
thresholds induce a containment relation for the sets of
matching queries. As we later rely on this property, we
state it as follows:

\begin{property}
\label{prop:subset}
Let $D$ be a stream database and $\theta_1$, $\theta_2$ be support thresholds. %with $\theta_1 \leq \theta_2$.
Let $Q_D^{\theta_1}$ and $Q_D^{\theta_2}$ be the sets of all matching queries. %, where $Q_i = \{q \mid \text{supp}(q,D) \geq \theta_i\}$ for $i \in \{1,2\}$.
If $\theta_2 \geq \theta_1$, then $Q_D^{\theta_2} \subseteq Q_D^{\theta_1}$.
\end{property}
% \textbf{Proposition:} Let $D$ be a stream database and $\theta_1$, $\theta_2$ be support thresholds. %with $\theta_1 \leq \theta_2$.
% Let $P_D^{\theta_1}$ and $P_D^{\theta_2}$ be the sets of all matching queries. %, where $Q_i = \{q \mid \text{supp}(q,D) \geq \theta_i\}$ for $i \in \{1,2\}$.
% If $\theta_2 \geq \theta_1$, then $P_D^{\theta_2} \subseteq P_D^{\theta_1}$. \\
\textit{Proof.}
    Let $q \in Q_D^{\theta_2}$. Then, supp$(q,D) \geq
    \theta_2 \geq \theta_1$, and thus, $q \in
    Q_D^{\theta_1}$
    and $Q_D^{\theta_2} \subseteq
    Q_D^{\theta_1}. $ \hfill \qed


\autoref{tab:example_stream_database} shows a database
of three streams from a cluster monitoring
application~\cite{reiss2012towards}. Each event has three
attributes, denoting the ID, status, and priority of a job
that is executed on a compute cluster. Now, consider the
query $q = \langle
 (x_0,\text{schedule},\text{low}),$
 $(\_,\text{schedule},\_),$ $(x_0,\text{evict},\text{low})
\rangle$. It matches streams $S_2$ and $S_3$, as, e.g.,
the three query terms of $q$ can be mapped to
events $e_{32}$, $e_{34}$, and $e_{35}$ of
$S_3$ (the mapped events are highlighted in blue). Note that
these are the only possible mappings for this scenario, since
the order of the query needs to be preserved.
Yet, the query does not match $S_1$, as the third query term
cannot be mapped to any event ($e_{14}$, highlighted
in red, does not qualify to be mapped).
With a support threshold of
0.6, query $q$ would be in the set of matching queries for
the stream database.

%\begin{wraptable}{r}{6cm}
\begin{table}
\centering
\small
\vspace{-.0em}
\caption{Example stream database.}
\label{tab:example_stream_database}
	\begin{tabular}{l lrrr}
		\toprule
		Stream & Event &  Job Id  &  Status  &  Priority\\
		\midrule
		$S_1$ & $e_{11}$   & \textbf{\color{cyan}\underline{1}} &
		\textbf{\color{cyan}\underline{schedule}} &
		\textbf{\color{cyan}\underline{low}}  \\
	 & $e_{12}$   & 1 & kill & low  \\
	 & $e_{13}$   & 1 & \textbf{\color{cyan}\underline{schedule}} &
	 high  \\
	 & $e_{14}$   & \textbf{\color{cyan}\underline{1}} &
	 \textbf{\color{red}\dotuline{finish}} &
	 \textbf{\color{red}\dotuline{high}}  \\
		\midrule
		$S_2$ & $e_{21}$   & \textbf{\color{cyan}\underline{2}} &
		\textbf{\color{cyan}\underline{schedule}} &
		\textbf{\color{cyan}\underline{low}}  \\
	 & $e_{22}$   & 3 & \textbf{\color{cyan}\underline{schedule}} &
	 high  \\
	 & $e_{23}$   & \textbf{\color{cyan}\underline{2}} &
	 \textbf{\color{cyan}\underline{evict}} & \textbf{\color{cyan}\underline{low}}
	 \\
		\midrule
		$S_3$ & $e_{31}$   & 4 & schedule & high  \\
	 & $e_{32}$   & \textbf{\color{cyan}\underline{5}} &
	 \textbf{\color{cyan}\underline{schedule}} &
	 \textbf{\color{cyan}\underline{low}}  \\
	 & $e_{33}$   & 4 & finish & high  \\
	 & $e_{34}$   & 6 & \textbf{\color{cyan}\underline{schedule}} &
	 low  \\
	 & $e_{35}$   & \textbf{\color{cyan}\underline{5}} &
	 \textbf{\color{cyan}\underline{evict}} &
	 \textbf{\color{cyan}\underline{low}}\\
		\bottomrule
\end{tabular}
%\vspace{-1em}
\end{table}




% \subsection{Query Discovery}
% \label{sec:discovery}

% Given a stream database, the problem of event query discovery relates to
% the identification of event queries that match to a sufficient number of streams to
% fulfill the support threshold.
% % , i.e., for
% % which there exists at least a single match for each stream.
% %are interested in queries for which there exists at least a single match of
% %the query for each stream.
% However, there may exist queries that fulfill the support threshold,
% but which are comparable in the sense that one of them
% is stricter than another one. Formally, a query $q$ is defined as stricter
% than a query $q'$, denoted by $q\prec q'$, if
% % $M(S,q)\subset M(S,q')$ for \emph{any} event stream $S$.
% %$\forall$ Stream $S:$ if $S \models q \Rightarrow S \models q', \exists \
% %S$ s.t. $S \models q$, but $S \not\models q$.
% (i) for any possible stream $S$ (not necessarily contained in a given stream
% database), $S \models q$ implies $S
% \models q'$, and (ii) there exists a stream $S$, such that $S \models q'$,
% but $S \not\models q$.

% In event query discovery, we are only interested in the strictest queries
% that fulfill the support threshold. The reasoning being that these queries
% denote the most concise characterization of the patterns of events
% that
% indicate a situation of interest. For a stream database $D$, we therefore
% consider a notion of descriptiveness~\cite{icdt2022}: A query $q$ is descriptive for $D$,
% if it is supported by $D$ and
% there does not exist a query $q'$ that is stricter, $q'\prec q$, and that is
% also supported by the stream database $D$.
% Based thereon, we formulate the problem of event query discovery:



% \begin{problem}[Event Query Discovery]
% \label{problem}
% Given a stream database $D=\{S_1,\ldots, S_d\}$ and a support threshold $\theta$,
% the problem of event query discovery is to construct the set of queries $Q_D^{\theta}$,
% such that:
% \begin{itemize}[nosep,left=1em]
%   \item $Q_D^{\theta}$ is correct: each $q\in Q_D^{\theta}$ creates at least a single match for
%   a sufficient number streams thus fulfilling the support threshold $\theta$;
%   , i.e., for all $q\in Q_D^{\theta}$ and $S\in M_+(q,D)$ it holds that $S \models
%   q$ and supp$(q,D) = \frac{|M_+(q,D)|}{|D|} \geq \theta$;
%   \item ${Q_D^{\theta}}$ is descriptive: only the strictest queries are considered,
%   i.e., for all $q,q'\in Q_D^{\theta}$ it holds that neither $q\prec q'$ nor $q'\prec q$;
%   \item ${Q_D^{\theta}}$ is complete: if a query $q$ is both correct and descriptive,
%   then it holds that $q\in Q_D^{\theta}$, i.e., ${Q_D^{\theta}Q}$ contains all queries
%   that are both correct and descriptive.
%   \end{itemize}
% \end{problem}







